\section{Introduction}
\label{sec:intro}

A 14-billion-parameter model in FP16 (16-bit floats) needs 29.5~GB for its
weights alone. That is more than any single consumer GPU holds. At batch 1,
serving one request at a time, a model reads every weight once for each token.
So the size of the weights sets both its memory and its speed. Two or three
bits per parameter let such a model run on local hardware.

Vector quantization rounds a group of weights to the nearest point of a fixed
set, the codebook. \citet{llvq2026} take that set from the Leech lattice
$\Lambda_{24}$, a regular arrangement of points in 24 dimensions. At two bits
per weight, their method gives the best quality in their comparison, ahead of
QuIP\#~\citep{quipsharp2024} and QTIP~\citep{qtip2024}. Each block of 24
weights becomes one 48-bit code, but the codebook has $1.1\times10^{14}$
points. Fast 2-bit GPU kernels decode with a small lookup table. A table that
large cannot be stored. In earlier work~\citep{llvq1preprint} we expanded each
code at load time into a wider format, which the kernel decodes with shifts
and masks. It worked, but the kernel read 4.804 bits per weight from GPU
memory (VRAM) for 2 bits of code. Most of the saving was lost on the GPU
(Figure~\ref{fig:gap}).

This paper removes that loss, then builds complete models around the result.

The first part is \Tetra{}, a new codebook on the same lattice. The lattice
is built from the Golay code, a binary code of length 24. Read in three groups
of eight coordinates, the code becomes a trellis: a layered graph with 64
nodes at each cut between groups, where each codeword is one path. A block is
then stored as an 8-bit state (the trellis state plus two parity bits), a few
bits that choose the path and the gain, and three 11-bit indices into one
16~KiB table. That is still 48 bits. The kernel decodes a block with six table
lookups and reads 2.148 bits per weight.

The second part builds complete model files. With the lattice codes
alone, Qwen3-4B scores about 15 points below FP16 on MMLU, a standard
multiple-choice test. Three changes close part of that gap without changing a
single code. We retrain the scale of each weight row to match the FP16 model.
We store the weight matrices (projections) that lose the most as 4-bit
integers. We store the embedding tables in 4 bits to make room for them. The
files stay near 2.7 bits per parameter over the whole model
(\S\ref{sec:file}).

\begin{figure}[t]
\centering
\includegraphics{fig_gap.pdf}
\caption{Bits of code a format carries per weight (hollow marker) and bits its
kernel reads per weight in VRAM (filled marker). Circles are our formats,
squares are deployed kernels. The three 2-bit formats carry 2.000 bits of
code. \Planes{}, our earlier layout, reads $2.40\times$ that because it
expands its codebook first. \Tetra{} reads $1.07\times$ that. The extra bits
are the row scales and a small tail of weights kept exactly in f16, not the
code. FP16 is left out: it carries and reads 16.000 bits per weight.
Data: echelle-formats.csv.}
\label{fig:gap}
\end{figure}

Our contributions:
\begin{itemize}
\item \Tetra{}, a codebook on $\Lambda_{24}$ read through a 64-state trellis
  of the Golay code and one 16~KiB table. We give the rule this design imposes
  on which points to keep, and an estimate of what it costs
  (\S\ref{sec:tetra}).
\item A fused kernel that decodes the weights and multiplies them by the
  input vector in one pass. It reads 2.148 bits per weight, against 4.804 for our earlier
  layout and 4.179 for 4-bit AWQ~\citep{awq2024} (\S\ref{sec:kernel}).
\item A recipe for complete files at about 2.7 bits per parameter. We measure
  each step on Qwen3-4B, 8B and 14B over the full MMLU test set. The
  confidence intervals are paired: each step and the one before it answer the
  same questions (\S\ref{sec:model}, \S\ref{sec:chain}).
\item Quality, speed and memory of our served files against FP16 and AWQ at
  all three sizes, and against 2-bit formats at 4B (\S\ref{sec:main}).
\end{itemize}

Our files do not match AWQ on quality. They score 4.76, 4.21 and 2.46 MMLU
points below it at 4B, 8B and 14B. The gap shrinks as the model grows. They
use about half its memory: 2.7 bits per parameter against 5.3 to 6.0. In our
own inference engine, they generate 113.8, 95.0 and 57.2 tokens per second
at 4B, 8B and 14B. Their weights take 1.38, 2.76 and 5.04~GB. At 4B they score 23.6 points
above IQ2\_XXS, a 2-bit llama.cpp format that uses 2.48 bits per parameter.
All measurements come from the same GPU, one NVIDIA L40S. For the experiments
behind the main results, we fixed the plan and our predictions before
measuring, and we report the predictions that were wrong
(Appendix~\ref{app:record}).
